\documentclass[12pt]{article}
\usepackage[spanish]{babel}
\usepackage[utf8]{inputenc}
\usepackage{graphicx}
\usepackage{geometry}
\usepackage{parskip}
\usepackage{setspace}
\usepackage{titlesec}
\usepackage{tabularx}
\usepackage{booktabs}
\usepackage{hyperref}
\usepackage{float}
\usepackage{longtable}
\usepackage{array}
\usepackage{colortbl}
\usepackage{xcolor}
\usepackage{enumitem}

\geometry{a4paper, margin=2.5cm}
\setlength{\parindent}{0pt}
\renewcommand{\baselinestretch}{1.15}

% Configuración de títulos
\titleformat{\section}{\large\bfseries}{}{0em}{}
\titleformat{\subsection}{\normalsize\bfseries}{}{0em}{}
\titleformat{\subsubsection}{\normalsize\bfseries\itshape}{}{0em}{}

% Definición de colores (valores entre 0 y 1)
\definecolor{high}{rgb}{0.9,0.2,0.2}
\definecolor{medium}{rgb}{0.95,0.7,0.2}
\definecolor{low}{rgb}{0.2,0.7,0.3}

% Para referencias con sangría francesa
\usepackage{quoting}
\newenvironment{hangparas}{\begin{quoting}[hang=true, indent=0.5cm]}{\end{quoting}}
\begin{document}
\setcounter{section}{1}
\setcounter{subsection}{4}
\subsection{Recomendación de modelo de madurez (CMMI / MoProSoft / ISO 33000)}

\subsubsection{Análisis comparativo de modelos}

A continuación se presenta una comparativa entre los tres modelos de madurez considerados para el contexto del proyecto Outline:

\begin{table}[H]
\centering
\small
\begin{tabularx}{\textwidth}{|l|X|X|X|}
\hline
\textbf{Criterio} & \textbf{CMMI-DEV v2.0} & \textbf{MoProSoft (NMX-I-059)} & \textbf{ISO/IEC 33000 (antigua 15504)} \\
\hline
Origen & EE.UU. (SEI/ISACA) & México (NYCE/UNAM) & Internacional (ISO) \\
\hline
Enfoque & Mejora de procesos para grandes contratistas (DoD) & PYMEs mexicanas de software & Evaluación de capacidad de procesos \\
\hline
Niveles & 5 niveles de madurez (0-5) & 3 niveles (Básico, Estándar, Avanzado) & 6 niveles de capacidad (0-5) \\
\hline
Costo de implementación & Alto (consultoría especializada) & Bajo a medio (documentación gratuita) & Medio (auditoría requerida) \\
\hline
Reconocimiento en México & Alto en empresas transnacionales & Alto en industria nacional (certificación NYCE) & Medio (usado para exportación) \\
\hline
Documentación requerida & Extensa (CMMI Practice Area) & Moderada (3 procesos: alta dirección, gerencia, operación) & Variable según objetivo \\
\hline
\end{tabularx}
\caption{Comparativa de modelos de madurez}
\end{table}

\subsubsection{Contexto específico de Outline}

El proyecto Outline (como producto open-source) y nuestro equipo (simulado de 5 estudiantes) tienen las siguientes características:

\begin{itemize}
    \item \textbf{Tamaño del equipo:} Muy pequeño (5 personas) - categoría PYME micro.
    \item \textbf{Sector:} Software de productividad, sin regulación específica (no aeroespacial, no médico).
    \item \textbf{Mercado objetivo:} México (empresa adoptante) + internacional (código abierto).
    \item \textbf{Ciclo de vida:} El equipo no desarrolla desde cero, sino que evalúa/adapta software existente.
\end{itemize}

\subsubsection{Recomendación: MoProSoft (con elementos de ISO 33000)}

\textbf{Justificación detallada:}

\begin{enumerate}
    \item \textbf{Alineación con el tamaño del equipo:} MoProSoft fue diseñado específicamente para PYMEs mexicanas de software (Oktaba \& Piattini, 2008). CMMI es excesivo para 5 personas (requiere roles como SEPG, QAG que no podemos cubrir).
    
    \item \textbf{Costo documental cero:} La norma NMX-I-059 está disponible públicamente en NYCE. Esto es crítico para un proyecto académico sin presupuesto.
    
    \item \textbf{Estructura simplificada:} MoProSoft organiza los procesos en tres categorías (Alta Dirección, Gerencia, Operación) que mapean directamente a:
    \begin{itemize}
        \item Alta Dirección $\rightarrow$ nuestro plan SQA (liderazgo del equipo)
        \item Gerencia $\rightarrow$ gestión de requisitos y planificación de proyectos
        \item Operación $\rightarrow$ desarrollo, pruebas y mantenimiento de Outline
    \end{itemize}
    
    \item \textbf{Relevancia para el mercado mexicano:} La empresa simulada que adoptaría Outline opera en Aguascalientes, México. Implementar prácticas de MoProSoft demuestra entendimiento del contexto nacional.
    
    \item \textbf{Complemento con ISO 33000:} Para la parte de medición de calidad, podemos tomar prestado el modelo de capacidad de procesos de ISO/IEC 33000 (niveles 0-5). Esto nos permite evaluar en qué nivel está nuestro proceso de pruebas, por ejemplo.
\end{enumerate}

\textbf{¿Por qué no CMMI?} Aunque CMMI es reconocido internacionalmente, requiere:

\begin{itemize}
    \item Un equipo dedicado a procesos (no factible con 5 personas).
    \item Entrenamiento formal (miles de dólares).
    \item Múltiples niveles de verificación (revisiones de pares, auditorías).
\end{itemize}

Esto no es práctico para un proyecto open-source mantenido por voluntarios.

\textbf{¿Por qué no solo ISO 33000?} ISO 33000 (antes 15504) es un modelo de evaluación más que de implementación. No prescribe prácticas específicas como MoProSoft; solo dice cómo medir capacidad. Necesitamos un modelo prescriptivo para guiar nuestras actividades diarias.

\subsubsection{Implementación propuesta de MoProSoft en nuestro contexto}

\begin{enumerate}
    \item \textbf{Procesos de Alta Dirección:}
    \begin{itemize}
        \item Gestión de calidad: seguimiento de KPIs (sección 2.3).
        \item Gestión de recursos: asignación de roles (tabla en \S 3 del plan SQA).
    \end{itemize}
    
    \item \textbf{Procesos de Gerencia:}
    \begin{itemize}
        \item Gestión de requisitos: backlog priorizado de HU (sección 2.1).
        \item Planificación de proyectos: cronograma de entregas (parcial 2 y 3).
    \end{itemize}
    
    \item \textbf{Procesos de Operación:}
    \begin{itemize}
        \item Desarrollo y mantenimiento: estándares de codificación (sección \S 5).
        \item Pruebas: estrategia definida en \S 7 del plan SQA.
    \end{itemize}
\end{enumerate}

En conclusión, MoProSoft complementado con ISO 33000 es la opción óptima para nuestro contexto por su adecuación a PYMEs, costo nulo, relevancia nacional y flexibilidad.

\subsection{Matriz de riesgos de calidad}

Se identifican 7 riesgos específicos del proyecto Outline, cada uno vinculado a una característica ISO 25010.

\begin{table}[H]
\centering
\footnotesize
\renewcommand{\arraystretch}{1.3}
\begin{tabularx}{\textwidth}{|l|X|c|c|X|}
\hline
\textbf{Riesgo} & \textbf{Característica afectada} & \textbf{Prob.} & \textbf{Impacto} & \textbf{Estrategia de mitigación} \\
\hline
R1: Configuración incorrecta de variables de entorno expone secretos (API keys, JWT secret) en logs o commits. & Seguridad & \cellcolor{medium}M & \cellcolor{high}A & Usar \texttt{git-secrets} pre-commit. Rotar secretos semanalmente. \\
\hline
R2: Editor colaborativo falla en conflictos de fusión, causando pérdida de datos. & Fiabilidad & \cellcolor{medium}M & \cellcolor{high}A & Pruebas de estrés con simulación de 10 usuarios concurrentes. Backups cada 5 min. \\
\hline
R3: Curva de aprendizaje alta para usuarios no técnicos. & Usabilidad & \cellcolor{high}A & \cellcolor{medium}M & Tutorial interactivo (tour guiado). Encuestas post-registro. \\
\hline
R4: Base de código de 50k+ líneas dificulta encontrar módulos específicos. & Mantenibilidad & \cellcolor{high}A & \cellcolor{medium}M & Generar documentación con Typedoc. Ownership de submódulos. \\
\hline
R5: Dependencias (npm) con vulnerabilidades críticas no detectadas. & Seguridad & \cellcolor{high}A & \cellcolor{high}A & Ejecutar \texttt{npm audit} y Snyk semanalmente. Política de actualización <7 días para CVSS >7. \\
\hline
R6: Rendimiento degradado en búsquedas cuando hay >10,000 documentos. & Eficiencia & \cellcolor{medium}M & \cellcolor{high}A & Indexar Elasticsearch. Pruebas de carga con k6. \\
\hline
R7: Falta de participación del equipo en actividades SQA. & Adecuación funcional & \cellcolor{high}A & \cellcolor{medium}M & Reunión semanal obligatoria de retrospectiva calidad (15 min). Checklist en cada PR. \\
\hline
\end{tabularx}
\caption{Matriz de riesgos de calidad para Outline}
\end{table}

\textbf{Leyenda:} \colorbox{high}{A = Alta} \quad \colorbox{medium}{M = Media} \quad \colorbox{low}{B = Baja}

\subsubsection{Plan de contingencia para riesgos críticos}

\begin{itemize}
    \item \textbf{R5 (vulnerabilidades en dependencias):} Si se detecta una vulnerabilidad crítica sin parche, se considera bloqueador de release y se busca workaround (ej. downgrade, parche manual).
    
    \item \textbf{R2 (pérdida de datos en colaboración):} Se implementa un \textit{shadow log} que registra todas las operaciones del editor; permite reconstruir documento desde cero.
\end{itemize}

% -- Fin de la parte que me tocó jaja... --

\end{document}
